\subsection{Endpoints}
\label{sec:endpoints}

A camada de rotas do módulo de solicitações foi implementada como um 
\textit{Blueprint} do Flask, registrado sob o prefixo \texttt{/solicitations}. 
A escolha do \textit{Blueprint} decorre diretamente da estratégia de 
\textit{Additive Change} descrita na Seção~\ref{sec:codigo-existente}: 
o módulo é incorporado como um componente isolado, sem alterar os 
controladores existentes. As rotas atuam como camada de fronteira — 
recebem a requisição, aplicam o controle de acesso, delegam o 
processamento ao \texttt{SolicitationService} e serializam a resposta —, 
mantendo fora do controlador as regras de negócio, a persistência e 
o tratamento de exceções de domínio. A documentação Swagger foi 
externalizada em arquivos YAML, referenciados pelo decorador 
\texttt{@swag\_from}, em contraste com o padrão observado nos 
controladores legados (Algoritmo~\ref{alg:legado}), nos quais a 
documentação ocupava a maior parte da função.

O Quadro~\ref{qua:endpoints} relaciona os endpoints implementados, 
a função responsável por cada um, a restrição de acesso aplicada e 
os requisitos funcionais atendidos.

\begin{quadro}[htbp]
\centering
\caption[Endpoints do módulo de solicitações]{Endpoints implementados no módulo de solicitações.}
\label{qua:endpoints}
\scriptsize
\setlength{\tabcolsep}{2pt}
\resizebox{\linewidth}{!}{%
\begin{tabular}{p{0.8cm}p{1.5cm}p{4.0cm}p{4.0cm}p{1.5cm}}
\toprule
\textbf{Método} & \textbf{Rota} & \textbf{Função} & \textbf{Controle de acesso} & \textbf{Requisitos} \\
\midrule
POST & \texttt{/} & \texttt{save\_solicitation} & Usuário autenticado & RF01, RF02, RF03 \\
GET & \texttt{/} & \texttt{get\_solicitations\_page} & Autenticado com perfil de visualização & RF04 \\
GET & \texttt{/me} & \texttt{get\_my\_solicitation} & Usuário autenticado (restrito ao próprio usuário) & RF04 \\
GET & \texttt{/\textless id\textgreater} & \texttt{get\_detail\_solicitation} & Autenticado, com verificação de acesso no serviço & RF04, RF11 \\
\bottomrule
\end{tabular}%
}
\fonte{Elaboração própria (2026). Obs: usando com báse que todas as rotas são prefixadas com \texttt{/solicitations}.}
\end{quadro}

A rota de criação (\texttt{POST /solicitations}) concentra o padrão 
adotado nas demais. O controlador limita-se a extrair os dados da 
requisição — aceitando tanto \texttt{multipart/form-data}, quando há 
arquivo anexado, quanto JSON, quando não há — e a delegar o 
processamento ao serviço, conforme ilustrado no Algoritmo~\ref{alg:rota-criacao}.

\begin{algorithm}[htbp]
\caption{Estrutura da rota de criação de solicitação.}
\label{alg:rota-criacao}
\begin{verbatim}
@bp.route("", methods=["POST"])
@auth.login_required
@swag_from("docs/create_solicitation.yml", validation=False)
def save_solicitation():
    form_data = request.form.to_dict() if request.form else None
    json_data = request.get_json(silent=True) or {}
    file = request.files.get("document_file")

    return SolicitationService.create_or_update(
        user_id=g.user.id,
        form_data=form_data,
        file=file,
    )
\end{verbatim}
\end{algorithm}

Quatro características do padrão merecem destaque. A primeira é o 
\textbf{controlador fino}: a função não constrói entidades, não valida 
campos individualmente, não executa consultas e não trata exceções de 
banco — responsabilidades que, no padrão legado (Algoritmo~\ref{alg:legado}), 
estavam concentradas na própria rota. A segunda é o \textbf{controle 
de acesso declarativo}, aplicado pelo decorador 
\texttt{@auth.login\_required}, que aceita um parâmetro opcional de 
perfil; a rota de listagem geral exige perfil de visualização, enquanto 
a rota \texttt{/me} restringe o resultado ao próprio usuário autenticado, 
independentemente do perfil. A terceira é a \textbf{validação na fronteira} 
por schemas Marshmallow: a consulta é interpretada por 
\texttt{SolicitationQuerySchema}, e a resposta é serializada por 
\texttt{SolicitationSchema} ou \texttt{PaginatedResponseSchema}, de modo 
que o formato de entrada e o de saída não dependem da interface que 
consome a API, atendendo ao RNF04. A quarta é a \textbf{externalização 
da documentação}: o decorador \texttt{@swag\_from} referencia um arquivo 
YAML por rota, com \texttt{validation=False} indicando que a validação 
é responsabilidade dos schemas e não do Swagger.

A listagem foi implementada por meio de uma função auxiliar, 
\texttt{\_fetch\_solicitations}, que centraliza a interpretação dos 
parâmetros de paginação, ordenação e filtro (situação, intervalo de 
datas e usuário). O endpoint \texttt{/solicitations/me} reutiliza a 
mesma função forçando o identificador do usuário autenticado, enquanto 
o endpoint \texttt{/solicitations} permite filtro opcional por usuário, 
desde que o solicitante possua o perfil de visualização. Essa 
centralização evita a duplicação da lógica de consulta entre rotas, 
que, no padrão legado, costumava ser reproduzida em cada controlador.

Do ponto de vista dos requisitos, os endpoints concentram os RF01 a RF04 
(criação e consulta) e o RF11 (restrição de operações a usuários 
autorizados). As operações de decisão, aplicação de efeitos, snapshot e 
reversão, correspondentes aos RF05 a RF09, são implementadas por rotas 
específicas descritas na Seção~\ref{sec:modulo-solicitacoes}.